% =====================================================================
%  Conditioning away ultra-small weights: repaired concentration lemma
%  (standalone LaTeX file you can paste into the paper)
% =====================================================================

\documentclass[11pt]{article}
\usepackage{amsmath,amssymb,amsthm,mathtools}
\usepackage{enumitem}

\newtheorem{thm}{Theorem}
\newtheorem{lem}{Lemma}
\newtheorem{prop}{Proposition}
\newtheorem{cor}{Corollary}
\newtheorem{rmk}{Remark}

\DeclareMathOperator{\E}{\mathbb E}
\DeclareMathOperator{\Pp}{\mathbb P}
\DeclareMathOperator{\Var}{\mathrm{Var}}

\begin{document}

\section*{Conditioning away ultra-small weights: repaired concentration}

\subsection*{Setup: exponential race and truncation sets}

Fix $\theta_1,\dots,\theta_n>0$ and let $E_1,\dots,E_n$ be independent exponentials
with rates $\theta_1,\dots,\theta_n$:
\[
\Pp(E_i>t)=e^{-\theta_i t},\qquad t\ge 0.
\]
Ordering $(E_i)$ increasingly produces the Luce permutation.
For $t\ge 0$ define the survivor count and survivor weight
\[
N(t):=\sum_{i=1}^n \mathbf 1\{E_i>t\},
\qquad
W(t):=\sum_{i=1}^n \theta_i\,\mathbf 1\{E_i>t\}.
\]

Fix two deterministic thresholds (in the intended application one takes
$\tau_n=(\log n)^{-2}$ and $L_n=(\log n)^2$):
\[
\tau_n\downarrow 0,\qquad L_n\uparrow\infty.
\]
Define the index sets
\[
S:=\{i:\theta_i<\tau_n\}\quad\text{(ultra-small)},\qquad
B:=\{i:\tau_n\le \theta_i\le L_n\}\quad\text{(bulk)},\qquad
H:=\{i:\theta_i>L_n\}\quad\text{(huge)}.
\]
Correspondingly, for $t\ge 0$ define
\[
N_B(t):=\sum_{i\in B}\mathbf 1\{E_i>t\},\qquad
W_B(t):=\sum_{i\in B}\theta_i\,\mathbf 1\{E_i>t\},
\]
and similarly $N_S,W_S$ and $N_H,W_H$.

\begin{rmk}[Why this is the right fix]
The bug in the naive ``$W(t)$ concentrates multiplicatively'' argument is that applying Bernstein
with the crude envelope $\sup_i \theta_i$ produces a denominator term of the form
``(envelope)$\times u$'' which dominates when $\E W(t)$ is small.
The point of restricting to $B$ is that \emph{both} a nontrivial lower bound
$\theta_i\ge \tau_n$ and an envelope $\theta_i\le L_n$ hold on $B$.
This restores concentration at the relevant quantile times on a scale governed by the
\emph{number of survivors} rather than the raw value of $\E W(t)$.
\end{rmk}

\subsection*{Step 1: conditioning does not hurt (it actually helps bookkeeping)}

Let $E_S:=(E_i)_{i\in S}$.
We will condition on $E_S$ (or even on $(E_i)_{i\in S\cup H}$).
The following observation is the formal version of your comment
``conditioning on small-$\theta$ only makes order-statistic locations \emph{more} concentrated.''

\begin{lem}[Independence after conditioning]\label{lem:cond-indep}
Condition on $E_S$ (and/or on $(E_i)_{i\in S\cup H}$).
Then the variables $(E_i)_{i\in B}$ remain independent exponentials with rates $(\theta_i)_{i\in B}$.
Consequently, any tail or concentration bound for $N_B(t)$ or $W_B(t)$ at deterministic times
holds \emph{verbatim} conditional on $E_S$.
\end{lem}

\begin{proof}
Independence of the family $(E_i)_{i=1}^n$ implies that conditioning on a subfamily does not
change the joint law of the complementary subfamily.
\end{proof}

\subsection*{Step 2: deterministic quantiles for the $B$-subsystem}

Define the deterministic mean curves
\begin{equation}\label{eq:HBDB}
H_B(t):=\frac{1}{n}\sum_{i\in B} e^{-\theta_i t},
\qquad
D_B(t):=\frac{1}{n}\sum_{i\in B} \theta_i e^{-\theta_i t}.
\end{equation}
Note that $H_B$ is decreasing, $H_B(0)=|B|/n$, and $H'_B(t)=-D_B(t)$.

For each integer $m\in\{1,\dots,|B|\}$ define the deterministic $B$-quantile time $t^B_m$ by
\begin{equation}\label{eq:tBm}
H_B(t^B_m)=\frac{m}{n}.
\end{equation}

Also define the \emph{random} $m$-survivor time in the $B$-subsystem:
\[
T^B_m:=\inf\{t\ge 0:\ N_B(t)\le m\}.
\]
(Equivalently, $T^B_m$ is the $(|B|-m)$-th order statistic of $(E_i)_{i\in B}$.)

\subsection*{Step 3: bracketing of $T^B_m$ by deterministic $t^B_{m\pm\Delta}$}

\begin{lem}[Quantile bracketing in $B$]\label{lem:bracket-B}
Fix $A>0$ and set $\Delta_m:=\lceil A\sqrt{m\log n}\rceil$.
Then for any sequence $M_n\le |B|$ satisfying $M_n\to\infty$,
\begin{equation}\label{eq:bracket-B}
\Pp\Big(\forall\, m\in[(\log n)^4,M_n]:\ t^B_{m+\Delta_m}\le T^B_m\le t^B_{m-\Delta_m}\Big)\to 1.
\end{equation}
Moreover, the same statement holds \emph{conditional on $E_S$} (and on $(E_i)_{i\in S\cup H}$).
\end{lem}

\begin{proof}
Fix $m$ and write $t^-:=t^B_{m+\Delta_m}$ and $t^+:=t^B_{m-\Delta_m}$.
By definition, $\E[N_B(t^-)] = nH_B(t^-)=m+\Delta_m$ and $\E[N_B(t^+)]=m-\Delta_m$.
Since $N_B(t)$ is a sum of independent Bernoullis, Chernoff bounds give
\[
\Pp(N_B(t^-)<m)\le \exp\!\Big(-c\frac{\Delta_m^2}{m+\Delta_m}\Big),
\qquad
\Pp(N_B(t^+)>m)\le \exp\!\Big(-c\frac{\Delta_m^2}{m}\Big),
\]
for a universal $c>0$.
With $\Delta_m\asymp \sqrt{m\log n}$, both probabilities are $\le n^{-c'A^2}$.
Choose $A$ so that $c'A^2>10$ and union bound over $m\le M_n\le n$ gives \eqref{eq:bracket-B}.
The conditional version follows from Lemma~\ref{lem:cond-indep}.
\end{proof}

\subsection*{Step 4: the repaired concentration lemma for $W_B(t)$}

This is the key replacement for the flawed ``relative concentration when the mean is tiny'' step.

\begin{lem}[Concentration of remaining weight in $B$ at $t^B_m$]\label{lem:WB-conc}
Fix $\eta\in(0,1/2)$.
There exists a constant $c_\eta>0$ such that for every $m\le |B|$,
\begin{equation}\label{eq:WB-conc-single}
\Pp\Big(\big|W_B(t^B_m) - nD_B(t^B_m)\big|\ge \eta\,nD_B(t^B_m)\Big)
\le
2\exp\!\Big(-c_\eta\,\frac{\tau_n}{L_n}\,m\Big).
\end{equation}
In particular, if $m\ge m_0(n):=C\frac{L_n}{\tau_n}\log n$ with $C$ large enough, then
\[
\Pp\Big(W_B(t^B_m)\in (1\pm\eta)\,nD_B(t^B_m)\Big)\ge 1-n^{-10}.
\]
Moreover, all statements above hold \emph{conditional on $E_S$} (and on $(E_i)_{i\in S\cup H}$).
\end{lem}

\begin{proof}
Fix $t\ge 0$ and write $W_B(t)=\sum_{i\in B}X_i(t)$ where
\[
X_i(t):=\theta_i\,\mathbf 1\{E_i>t\}.
\]
Then $X_i(t)\in[0,\theta_i]\subset[0,L_n]$, and the $X_i(t)$ are independent
(and remain so conditional on $E_S$ by Lemma~\ref{lem:cond-indep}).
Moreover
\[
\mu(t):=\E[W_B(t)] = \sum_{i\in B}\theta_i e^{-\theta_i t}=nD_B(t).
\]

We apply Bernstein's inequality to $\sum_{i\in B}X_i(t)$.
First bound the variance proxy:
\[
v(t):=\sum_{i\in B}\Var(X_i(t))
\le \sum_{i\in B}\E[X_i(t)^2]
=\sum_{i\in B}\theta_i^2 e^{-\theta_i t}
\le L_n\sum_{i\in B}\theta_i e^{-\theta_i t}
= L_n\,\mu(t),
\]
since $\theta_i^2\le L_n\theta_i$ on $B$.

Bernstein (with envelope $L_n$) gives for $u>0$:
\[
\Pp\big(|W_B(t)-\mu(t)|\ge u\big)
\le 2\exp\!\Big(-\frac{u^2}{2(v(t)+L_n u/3)}\Big).
\]
Set $t=t^B_m$ and $u=\eta\mu(t^B_m)=\eta nD_B(t^B_m)$. Using $v(t^B_m)\le L_n\mu(t^B_m)$ we obtain
\[
\Pp\Big(|W_B(t^B_m)-nD_B(t^B_m)|\ge \eta nD_B(t^B_m)\Big)
\le
2\exp\!\Big(-c\,\eta^2\frac{\mu(t^B_m)}{L_n}\Big)
\]
for a universal $c>0$.

Finally, use the crucial \emph{lower bound on the mean} coming from the truncation:
\[
\mu(t^B_m)=\sum_{i\in B}\theta_i e^{-\theta_i t^B_m}
\ge
\tau_n\sum_{i\in B} e^{-\theta_i t^B_m}
=\tau_n\,nH_B(t^B_m)
=\tau_n\,m,
\]
where we used $\theta_i\ge\tau_n$ on $B$ and the definition \eqref{eq:tBm}.
Thus the exponent is at least $c_\eta(\tau_n/L_n)m$, proving \eqref{eq:WB-conc-single}.
The conditional version is identical by Lemma~\ref{lem:cond-indep}.
\end{proof}

\subsection*{Step 5: uniform concentration and the ``restored'' regime}

\begin{cor}[Uniform-in-$m$ concentration on a growing range]\label{cor:WB-uniform}
Fix $\eta\in(0,1/2)$ and set $m_0(n):=C\frac{L_n}{\tau_n}\log n$ for $C$ large enough.
Then for any $M_n\le |B|$,
\[
\Pp\Big(\forall\, m\in[m_0(n),M_n]:\ W_B(t^B_m)\in(1\pm\eta)\,nD_B(t^B_m)\Big)\to 1.
\]
The same holds conditional on $E_S$ (and on $(E_i)_{i\in S\cup H}$).
\end{cor}

\begin{proof}
Apply Lemma~\ref{lem:WB-conc} and union bound:
\[
\sum_{m=m_0(n)}^{M_n}2\exp\!\Big(-c_\eta\frac{\tau_n}{L_n}m\Big)
\le
\frac{2\exp\!\big(-c_\eta\frac{\tau_n}{L_n}m_0(n)\big)}{1-\exp(-c_\eta\tau_n/L_n)}
\le n^{-9}
\]
for $C$ sufficiently large. The conditional statement follows from Lemma~\ref{lem:cond-indep}.
\end{proof}

\begin{rmk}[How this plugs back into the fixed-point method]
This is exactly the regime needed in the ``zoom at $t_m$'' method:
once you restrict to the $B$-subsystem (after conditioning on $E_S$), the denominator
that previously failed to concentrate now does concentrate uniformly for all
$m\ge m_0(n)=C(L_n/\tau_n)\log n$.
The remaining finitely many $m<m_0(n)$ correspond to only polylogarithmically many
end-steps and can be handled by crude bounds (or absorbed into an $o(1)$ term if you
are working at a scale $\gg \log n$).
\end{rmk}

\end{document}